\documentclass[11pt,letterpaper]{article}
\usepackage[T1]{fontenc}
\usepackage[utf8]{inputenc}
\usepackage[scaled=0.95]{helvet}
\renewcommand{\familydefault}{\sfdefault}
\usepackage[top=1.7cm,bottom=1.5cm,left=2.2cm,right=2.2cm]{geometry}
\usepackage{microtype}
\usepackage[hidelinks]{hyperref}
\pagestyle{empty}
\setlength{\parindent}{0pt}
\setlength{\parskip}{7pt}

\begin{document}

{\large\bfseries Grover Aruquipa}\par\vspace{-6pt}
{\small Qu\'ebec, QC, Canada \ \textbullet\ \ grover.aruquipa.9@gmail.com \ \textbullet\ \ groveraruquipa.github.io}

\vspace{4pt}
Hiring Team, Intuitive, Strategy and Growth (Future Forward)\\
Mechanical Engineer, Robotic Accessories

Dear Hiring Team,

For the last three years I have been working on one mechanical problem: how to roll a tool without limit while still opening and closing a jaw at its tip, with every motor kept away from the tool. A wristed instrument on a da Vinci system solves a version of that problem in every procedure. Designing the accessories and mechanisms around it is the work I want to do next.

I am finishing my PhD at Universit\'e Laval with Prof.\ Cl\'ement Gosselin and Prof.\ Patrice Lambert, and I defend in early 2027; I also hold an MSc from France. I led the full development of three robotic mechanisms with 4, 7 and 8 actuators. For each one I wrote the requirements, defined the architecture, did the SolidWorks design and the stress analysis, selected the motors, encoders, gearboxes and bearings, machined and printed the parts, assembled them, built the test setups and ran the verification tests, then redesigned what the tests showed was weak.

The mechanism closest to your products is a 4-DoF coaxial wrist-gripper. It sends four actuation channels through concentric hollow shafts in a 180\,mm footprint, with 0.25\,kg of moving mass and a non-backdrivable jaw, and it turned its tool ten full revolutions with less than one degree of tracking error while holding the object. Another of my robots actuates its gripper from the base through a concentric tube while the end effector rotates without limit. I am now building a Delta robot with a Sarrus gripper that estimates grasp force from the motor currents, without a sensor at the jaw, which is the kind of capability that matters when the tool is touching tissue.

I am used to documenting a design well enough for someone else to judge it. Every one of my designs has been reviewed before fabrication by my supervisors and after by referees of the ASME Journal of Mechanisms and Robotics and Mechanism and Machine Theory, and I review for that journal myself. I also mentor two graduate students day to day on an 8-actuator robot, so I know how to guide other engineers through a design and how to sit on the other side of a review.

I am available from January 2027 and ready to relocate to Sunnyvale and work fully on site. Prof.\ Gosselin, Prof.\ Lambert and Univ.-Prof.\ Andreas M\"uller of JKU Linz can speak to my work; their details are in my CV.

Thank you for your time. I would be glad to show you the prototypes on video, or to bring the wrist-gripper to an interview.

Sincerely,\\
Grover Aruquipa

\end{document}
